% M6共存試験: hyperref + amsmath の \eqref（和文コンテキスト中の式番号参照）。
% 検証するもの:
%   (a) 和文に挟まれた位置での \eqref/\ref（ASCIIラベル名）の解決
%   (b) \text{和文} を含む式と、直打ち和文を含む式（backend M1で後者も組める）
%   (c) 和文見出し（bookmarks 同様 UTF-16BE 経路）と式番号の共存
% 既知の制限（ワークアラウンドを本コメントで示す）:
%   * ラベル「名」への和文使用（\label{eq:あ}）は不可。CJKのアクティブ文字が
%     LaTeXカーネルの \csname r@<name>\endcsname 内で展開され，バイトマクロの
%     内部で「Missing \endcsname inserted」となる。CJKutf8 共通でbackend非依存。
%     （(u)pTeX ではバイト非アクティブのため通る。）
%     ワークアラウンド: ラベル名はASCIIにする（本fixtureのとおり）。本文・式・
%     見出しの和文はすべて使える。
\documentclass[pdflatex]{jlreq}
\usepackage{amsmath}
\usepackage[unicode]{hyperref}
\begin{document}
\section{和文の見出しと式番号}\label{sec:top}
和文に挟まれた参照\eqref{eq:euler}と、節の参照\ref{sec:top}。
\begin{equation}
e^{i\pi}+1=0 \label{eq:euler}
\end{equation}
和文を含む式の参照\eqref{eq:distance}:
\begin{equation}
\text{道のり} = \text{速さ} \times \text{時間} \label{eq:distance}
\end{equation}
直打ち和文を含む式の参照\eqref{eq:einstein}:
\begin{equation}
E = mc^2 \quad アインシュタイン \label{eq:einstein}
\end{equation}
式\eqref{eq:euler}、\eqref{eq:distance}、\eqref{eq:einstein}をまとめて参照。
\end{document}
